\documentclass[11pt]{amsart}

\usepackage[a4paper,margin=1in]{geometry}
\usepackage{amsmath,amssymb,amsthm,mathtools}
\usepackage{booktabs}
\usepackage{enumitem}
\usepackage[hidelinks]{hyperref}

\title{An Ordinal-Prime Transform of the M\"obius Function}
\author{}
\date{}

\newtheorem{definition}{Definition}[section]
\newtheorem{proposition}[definition]{Proposition}
\newtheorem{corollary}[definition]{Corollary}
\newtheorem{remark}[definition]{Remark}

\newcommand{\mob}{\mu}
\newcommand{\Mert}{M}
\newcommand{\Sq}{Q}
\newcommand{\PS}{\mathcal P}

\begin{document}

\maketitle
\tableofcontents

\section{Purpose and scope}

This note develops only the ordinal-prime transform introduced in the present line of discussion. It is independent of the earlier Leslie-matrix constructions and does not use them.

The starting point is the M\"obius function
\[
\mob(n)\in\{-1,0,+1\}.
\]
The occurrences of $-1$ and $+1$ are numbered separately, in chronological order, and the corresponding ordinal prime is attached to each occurrence. Zeros are replaced by $1$.

The resulting signed sequence is denoted by $S(n)$, and its cumulative sum by
\[
SS(n)=\sum_{k\le n}S(k).
\]

The main conclusions are:

\begin{enumerate}[label=(\roman*)]
\item $S(n)$ is an ordinal-prime reweighting of the three-valued M\"obius sequence.
\item $SS(n)$ has an exact closed expression in terms of the Mertens function, the squarefree counting function, and the sum of the first $m$ primes.
\item After removing the zero contribution, $SS(n)$ is exactly a signed sum of $|M(n)|$ consecutive primes.
\item Consequently, $SS(n)$ is essentially an amplified nonlinear reparametrization of $M(n)$ once $Q(n)$ is known.
\item The transform does not create independent arithmetic information by itself; any usefulness must come from a new structural property of the transformed representation.
\end{enumerate}

\section{Definition of the three rows}

Let
\[
p_1=2,\qquad p_2=3,\qquad p_3=5,\qquad p_4=7,\ldots
\]
denote the primes in increasing order.

Define the cumulative occurrence counts
\[
N_-(n)
=
\sum_{k\le n}\frac{\mob(k)^2-\mob(k)}2,
\]
\[
N_0(n)
=
\sum_{k\le n}\bigl(1-\mob(k)^2\bigr),
\]
and
\[
N_+(n)
=
\sum_{k\le n}\frac{\mob(k)^2+\mob(k)}2.
\]

Thus $N_-(n)$ counts the occurrences of $\mob=-1$ up to $n$, $N_0(n)$ counts the zeros, and $N_+(n)$ counts the occurrences of $\mob=+1$.

The three rows are

\[
R_-(n)
=
-\frac{\mob(n)^2-\mob(n)}2\,p_{N_-(n)},
\]

\[
R_0(n)
=
1-\mob(n)^2,
\]

and

\[
R_+(n)
=
\frac{\mob(n)^2+\mob(n)}2\,p_{N_+(n)}.
\]

Equivalently,

\[
R_-(n)=
\begin{cases}
-p_{N_-(n)},&\mob(n)=-1,\\
0,&\text{otherwise},
\end{cases}
\]

\[
R_0(n)=
\begin{cases}
1,&\mob(n)=0,\\
0,&\text{otherwise},
\end{cases}
\]

and

\[
R_+(n)=
\begin{cases}
p_{N_+(n)},&\mob(n)=+1,\\
0,&\text{otherwise}.
\end{cases}
\]

\section{Definition of the ordinal-prime transform}

\begin{definition}
The ordinal-prime transform $S(n)$ is
\[
S(n):=R_-(n)+R_0(n)+R_+(n).
\]
Hence
\[
\boxed{
S(n)=
\begin{cases}
-p_{N_-(n)},&\mob(n)=-1,\\[1mm]
1,&\mob(n)=0,\\[1mm]
+p_{N_+(n)},&\mob(n)=+1.
\end{cases}
}
\]
\end{definition}

The first occurrences therefore behave as follows:

\[
\mob(n)=-1
\quad\Longrightarrow\quad
-2,-3,-5,-7,-11,\ldots
\]

in ordinal order, whereas

\[
\mob(n)=+1
\quad\Longrightarrow\quad
2,3,5,7,11,\ldots
\]

in its own separate ordinal order.

Zeros are mapped to $1$.

\section{The cumulative transform}

\begin{definition}
Define
\[
\boxed{
SS(n):=\sum_{k=1}^{n}S(k).
}
\]
\end{definition}

Let
\[
\PS(m):=\sum_{j=1}^{m}p_j
\]
be the sum of the first $m$ primes.

Then, by construction,

\begin{proposition}
For every $n\ge1$,
\[
\boxed{
SS(n)
=
\PS\bigl(N_+(n)\bigr)
-
\PS\bigl(N_-(n)\bigr)
+
N_0(n).
}
\]
\end{proposition}

\begin{proof}
Each positive M\"obius occurrence contributes the next unused positive prime, so the total positive contribution is $\PS(N_+(n))$. Each negative occurrence contributes the negative of the next unused prime, giving $-\PS(N_-(n))$. Every zero contributes $1$, producing $N_0(n)$.
\end{proof}

\section{Elimination of the occurrence counters}

Introduce the standard functions
\[
\Mert(n):=\sum_{k\le n}\mob(k)
\]
and
\[
\Sq(n):=\sum_{k\le n}\mob(k)^2.
\]

Since
\[
N_+(n)+N_-(n)=\Sq(n)
\]
and
\[
N_+(n)-N_-(n)=\Mert(n),
\]
one has the exact identities
\[
\boxed{
N_+(n)=\frac{\Sq(n)+\Mert(n)}2,
\qquad
N_-(n)=\frac{\Sq(n)-\Mert(n)}2.
}
\]

Also,
\[
\boxed{
N_0(n)=n-\Sq(n).
}
\]

Therefore:

\begin{proposition}
The cumulative ordinal-prime transform satisfies
\[
\boxed{
SS(n)
=
\PS\!\left(\frac{\Sq(n)+\Mert(n)}2\right)
-
\PS\!\left(\frac{\Sq(n)-\Mert(n)}2\right)
+
n-\Sq(n).
}
\]
\end{proposition}

This is the central exact identity.

If one wishes to avoid all auxiliary notation except $\mob$ and the prime sequence, it can also be written as

\[
\boxed{
\begin{aligned}
SS(n)
={}&
\sum_{j=1}^{\frac12\sum_{k\le n}(\mob(k)^2+\mob(k))}p_j
\\
&-
\sum_{j=1}^{\frac12\sum_{k\le n}(\mob(k)^2-\mob(k))}p_j
\\
&+
\sum_{k\le n}\bigl(1-\mob(k)^2\bigr).
\end{aligned}
}
\]

\section{Signed block of consecutive primes}

The exact identity becomes particularly transparent after removing the zero contribution. Define

\[
T(n):=SS(n)-\bigl(n-\Sq(n)\bigr).
\]

Then

\[
T(n)
=
\PS\!\left(\frac{\Sq(n)+\Mert(n)}2\right)
-
\PS\!\left(\frac{\Sq(n)-\Mert(n)}2\right).
\]

\begin{proposition}
If $\Mert(n)>0$, then
\[
\boxed{
T(n)
=
\sum_{j=(\Sq(n)-\Mert(n))/2+1}^{(\Sq(n)+\Mert(n))/2}p_j.
}
\]

If $\Mert(n)<0$, then
\[
\boxed{
T(n)
=
-
\sum_{j=(\Sq(n)-|\Mert(n)|)/2+1}^{(\Sq(n)+|\Mert(n)|)/2}p_j.
}
\]

If $\Mert(n)=0$, then $T(n)=0$.
\end{proposition}

Hence $T(n)$ is exactly the signed sum of $|\Mert(n)|$ consecutive primes, centered at index approximately $\Sq(n)/2$.

Two immediate consequences are

\[
\boxed{
\operatorname{sgn}T(n)=\operatorname{sgn}\Mert(n)
}
\]

and

\[
\boxed{
T(n)=0\iff \Mert(n)=0.
}
\]

Thus the transform preserves exactly the sign and the zero set of the Mertens function after subtraction of the deterministic zero-count contribution.

\section{Exact prime-block bounds}

Because the block contains exactly $|\Mert(n)|$ consecutive primes, one has

\[
\boxed{
|\Mert(n)|\,
p_{\frac{\Sq(n)-|\Mert(n)|}{2}+1}
\le
|T(n)|
\le
|\Mert(n)|\,
p_{\frac{\Sq(n)+|\Mert(n)|}{2}}.
}
\]

These inequalities are completely elementary but useful: they show that bounds for $T(n)$ and bounds for $\Mert(n)$ are directly interconvertible once the prime scale around $\Sq(n)/2$ is controlled.

\section{Asymptotic scale}

The squarefree counting function satisfies

\[
\Sq(n)\sim \frac{6}{\pi^2}n.
\]

Hence

\[
\frac{\Sq(n)}2
\sim
\frac{3}{\pi^2}n.
\]

Also, by the prime number theorem,

\[
p_m\sim m\log m.
\]

Therefore the characteristic prime scale of the block defining $T(n)$ is

\[
p_{\Sq(n)/2}
\sim
\frac{3}{\pi^2}n\log n.
\]

Since $\Mert(n)=o(n)$, the block of $|\Mert(n)|$ primes is short relative to its center on the index scale. Consequently,

\[
\boxed{
T(n)\sim \Mert(n)\,p_{\Sq(n)/2}
}
\]

in the natural first-order sense, and thus

\[
\boxed{
SS(n)-n+\Sq(n)
\sim
\frac{3}{\pi^2}n\log n\,\Mert(n).
}
\]

Accordingly, the raw cumulative transform behaves as an amplified version of the Mertens function, with amplification factor of order $n\log n$.

\section{Big-O relation}

The exact block representation immediately gives

\[
|T(n)|
=
O\!\bigl(n\log n\,|\Mert(n)|\bigr).
\]

Hence

\[
\boxed{
SS(n)
=
O\!\bigl(n\log n\,|\Mert(n)|+n\bigr).
}
\]

Using only the trivial estimate $|\Mert(n)|\le n$ gives

\[
SS(n)=O(n^2\log n).
\]

Using $\Mert(n)=o(n)$ gives

\[
SS(n)=o(n^2\log n).
\]

\section{Conditional scale under the Riemann hypothesis}

Assuming the Riemann hypothesis, one has for every $\varepsilon>0$,

\[
\Mert(n)=O_\varepsilon\!\left(n^{1/2+\varepsilon}\right).
\]

Therefore

\[
SS(n)
=
O_\varepsilon\!\left(n^{3/2+\varepsilon}\log n\right).
\]

Absorbing the logarithm into $n^\varepsilon$,

\[
\boxed{
SS(n)
=
O_\varepsilon\!\left(n^{3/2+\varepsilon}\right)
\qquad
(\forall\,\varepsilon>0).
}
\]

It is important not to replace this by the stronger statement
\[
\Mert(n)=O(\sqrt n),
\]
which is not implied by the Riemann hypothesis.

\section{A corresponding RH reformulation}

Because
\[
T(n)=SS(n)-n+\Sq(n)
\]
is asymptotically $\Mert(n)$ multiplied by a prime of size $\asymp n\log n$, the Mertens criterion for RH can be transported to the transformed sequence.

At the natural $\varepsilon$-level, one obtains the equivalent scale

\[
\boxed{
\mathrm{RH}
\Longleftrightarrow
SS(n)-n+\Sq(n)
=
O_\varepsilon\!\left(n^{3/2+\varepsilon}\right)
\quad
(\forall\,\varepsilon>0),
}
\]
with the usual understanding that logarithmic factors are absorbed into $\varepsilon$.

This is a reformulation, not an independent proof mechanism.

\section{Insertion of the Riemann zeta function}

The two M\"obius-derived ingredients of the exact identity have standard Dirichlet-series representations:

\[
\sum_{m\ge1}\frac{\mob(m)}{m^s}
=
\frac1{\zeta(s)},
\qquad
\Re s>1,
\]

and

\[
\sum_{m\ge1}\frac{\mob(m)^2}{m^s}
=
\frac{\zeta(s)}{\zeta(2s)},
\qquad
\Re s>1.
\]

Perron's formula therefore gives, formally and with the standard truncation/endpoint qualifications,

\[
\Mert(x)
=
\frac{1}{2\pi i}
\int_{c-i\infty}^{c+i\infty}
\frac{x^s}{s\zeta(s)}\,ds,
\qquad
c>1,
\]

and

\[
\Sq(x)
=
\frac{1}{2\pi i}
\int_{c-i\infty}^{c+i\infty}
\frac{\zeta(s)}{\zeta(2s)}
\frac{x^s}{s}\,ds.
\]

Substituting these into the exact identity yields a zeta-functional representation of $SS$:

\[
\begin{aligned}
SS(n)
={}&
\PS\!\left[
\frac{1}{4\pi i}
\int_{c-i\infty}^{c+i\infty}
\left(
\frac{\zeta(s)}{\zeta(2s)}
+
\frac1{\zeta(s)}
\right)
\frac{n^s}{s}\,ds
\right]
\\
&-
\PS\!\left[
\frac{1}{4\pi i}
\int_{c-i\infty}^{c+i\infty}
\left(
\frac{\zeta(s)}{\zeta(2s)}
-
\frac1{\zeta(s)}
\right)
\frac{n^s}{s}\,ds
\right]
\\
&+
n
-
\frac{1}{2\pi i}
\int_{c-i\infty}^{c+i\infty}
\frac{\zeta(s)}{\zeta(2s)}
\frac{n^s}{s}\,ds.
\end{aligned}
\]

The prime-sum function $\PS(m)$ can itself be represented through the prime-counting function, and hence through $\log\zeta(s)$. One convenient identity is

\[
\PS(m)
=
m\,p_m
-
\int_2^{p_m}\pi(t)\,dt,
\]

with the ordinary discrete interpretation at the jumps of $\pi$.

The prime-counting structure can in turn be recovered from the classical prime-power counting function built from $\log\zeta$.

Thus three familiar zeta structures enter:

\[
\boxed{\frac1{\zeta(s)}}
\]

for the signed M\"obius imbalance,

\[
\boxed{\frac{\zeta(s)}{\zeta(2s)}}
\]

for squarefree counting, and

\[
\boxed{\log\zeta(s)}
\]

for the prime weights.

\section{Are the three zeta structures themselves new?}

No. Each component is classical:

\begin{enumerate}[label=(\alph*)]
\item $1/\zeta(s)$ is the Dirichlet series of $\mob$.
\item $\zeta(s)/\zeta(2s)$ is the Dirichlet series of $\mob^2$.
\item $\log\zeta(s)$ encodes primes and prime powers.
\end{enumerate}

What is specific to the present construction is the nonlinear composition

\[
\PS\!\left(\frac{\Sq+\Mert}{2}\right)
-
\PS\!\left(\frac{\Sq-\Mert}{2}\right)
+
n-\Sq.
\]

However, the fact that the combination is nonstandard does not by itself imply that it contains new arithmetic information.

\section{Why mixing the characterized pieces is not automatically useful}

Once $n$ is fixed, the pair
\[
(\Mert(n),\Sq(n))
\]
already determines the three occurrence counts:

\[
N_+(n)=\frac{\Sq(n)+\Mert(n)}2,
\]

\[
N_-(n)=\frac{\Sq(n)-\Mert(n)}2,
\]

and

\[
N_0(n)=n-\Sq(n).
\]

Therefore $SS(n)$ is a deterministic transformation of $(\Mert(n),\Sq(n))$:

\[
\boxed{
(\Mert(n),\Sq(n))
\longmapsto
SS(n).
}
\]

In particular, once $\Sq(n)$ is known, the centered transform

\[
T(n)=SS(n)-n+\Sq(n)
\]
is strictly monotone in $\Mert(n)$.

Thus the construction does not provide an independent arithmetic constraint on $\Mert(n)$. It reweights and amplifies the same imbalance.

This point is essential. Merely combining two completely characterized objects does not create a new source of information. For a transformed representation to become genuinely useful, it would need to produce an additional structural property such as:

\begin{enumerate}[label=(\roman*)]
\item a new cancellation unavailable in the original variables;
\item a positivity or monotonicity principle strong enough to yield new bounds;
\item a spectral representation in which a difficult Mertens bound becomes a simpler geometric or operator inequality;
\item an inequality between independently controlled components;
\item or a transformation in which the prime side constrains the M\"obius side in a way not already encoded by $(\Mert,\Sq)$.
\end{enumerate}

No such independent mechanism has yet been extracted from $SS(n)$ itself.

\section{A normalized form}

A natural normalization is

\[
\boxed{
R(n)
:=
\frac{SS(n)-n+\Sq(n)}
{p_{\lfloor \Sq(n)/2\rfloor}}.
}
\]

Since
\[
SS(n)-n+\Sq(n)
\sim
\Mert(n)\,p_{\Sq(n)/2},
\]
one expects

\[
\boxed{
R(n)\sim \Mert(n).
}
\]

This normalization removes the dominant $n\log n$ prime amplification.

Again, this is primarily a reparametrization of Mertens, not an independent object.

\section{A prime-curvature residual}

One can isolate the deviation from replacing the entire consecutive-prime block by a constant central prime. Define

\[
\boxed{
E(n)
:=
SS(n)-n+\Sq(n)
-
\Mert(n)p_{\lfloor \Sq(n)/2\rfloor}.
}
\]

Then $E(n)$ measures the local variation of the primes across a window of length $|\Mert(n)|$ centered near index $\Sq(n)/2$.

Indeed, when $\Mert(n)>0$,

\[
E(n)
=
\sum_{j=(\Sq-\Mert)/2+1}^{(\Sq+\Mert)/2}
\left(
p_j-p_{\lfloor \Sq/2\rfloor}
\right),
\]

and for negative $\Mert(n)$ the corresponding signed formula holds.

This residual separates two ingredients:

\[
\boxed{\text{Mertens window length}}
\qquad\text{and}\qquad
\boxed{\text{local prime variation}}.
\]

Whether this residual admits useful bounds that are not merely restatements of known information is an open direction within this construction.

\section{Information-theoretic interpretation}

At each fixed $n$, $SS(n)$ does not retain the internal ordering of the previous $+1$ and $-1$ values of the M\"obius sequence.

Only the totals matter:

\[
N_+(n),\qquad N_-(n),\qquad N_0(n).
\]

Consequently, two M\"obius sign histories with identical endpoint counts produce the same $SS(n)$.

Therefore the cumulative transform is not sensitive to temporal ordering within the M\"obius sequence. It is a nonlinear endpoint statistic.

This explains both its simplicity and its limitation.

\section{Novelty status}

The standard ingredients used above are classical. The particular ordinal-prime assignment

\[
-1\mapsto -p_{\text{ordinal among }-1\text{ occurrences}},
\]

\[
0\mapsto1,
\]

\[
+1\mapsto p_{\text{ordinal among }+1\text{ occurrences}}
\]

and the resulting sequence $SS(n)$ were not identified in the preceding search as a standard named sequence or classical transform.

This should be treated as a novelty observation about the specific formulation, not as a claim that its consequences are automatically new. The exact reduction to $(M,Q)$ shows that many of its first-order properties follow directly from classical objects.

\section{Summary}

The construction can be summarized by the exact chain

\[
\mob
\longrightarrow
(N_-,N_0,N_+)
\longrightarrow
S(n)
\longrightarrow
SS(n),
\]

with

\[
N_\pm(n)=\frac{\Sq(n)\pm\Mert(n)}2,
\qquad
N_0(n)=n-\Sq(n),
\]

and

\[
\boxed{
SS(n)
=
\PS\!\left(\frac{\Sq(n)+\Mert(n)}2\right)
-
\PS\!\left(\frac{\Sq(n)-\Mert(n)}2\right)
+n-\Sq(n).
}
\]

After removing the zero contribution,

\[
\boxed{
SS(n)-n+\Sq(n)
=
\text{signed sum of exactly }|\Mert(n)|\text{ consecutive primes}.
}
\]

Hence

\[
\boxed{
SS(n)-n+\Sq(n)
\sim
\frac{3}{\pi^2}n\log n\,\Mert(n).
}
\]

The transform therefore amplifies the Mertens imbalance by the local prime scale, but does not by itself introduce an independent constraint capable of improving known Mertens bounds.

The most structurally distinct residual produced so far is

\[
\boxed{
E(n)
=
SS(n)-n+\Sq(n)
-
\Mert(n)p_{\lfloor \Sq(n)/2\rfloor},
}
\]

which measures the interaction between the Mertens window length and local variation of consecutive primes.

\end{document}
